\begin{abstract}

The SuperKEKB accelerator is an electron-positron collider that has achieved world record instantaneous luminosity of $5.3 \times 10^{34}cm^{-2}s^{-1}$. Injection beam loss is the one of the challenges in increasing the beam currents and achieving the design luminosity. However, its underlying mechanisms are not yet fully understood.

In this study, we developed a monitoring system dedicated to injection beam loss by installing sensors at multiple collimators around the SuperKEKB rings. The system operated continuously throughout the experimental period and was used to support injection and collimator tuning. Using this system, we measured beam loss as a function of collimator position to investigate the collimator offset. This enabled to get offset value at the collimator which has a difficulty to get the value. 

We also performed frequency analysis of injection beam loss signals to investigate their time structure. The analysis revealed a frequency component corresponding to the horizontal betatron tune in the vertical injection beam loss signals. This behavior was also observed in simulations, providing further insight into the dynamics of the injected beam. 

\end{abstract}